\section{Conclusion and Future Work}

% We introduced Parametric Memory Distillation, a retrieval to memory paradigm that compiles the benefits of retrieval augmentation into a lightweight, learnable module for adapting frozen Time Series Foundation Models.
% Our instantiation, TS-Memory, distills both retrieval-induced uncertainty and local domain regularities from a leakage-safe $k$NN teacher into a plug-and-play memory module via confidence-gated supervision.
% Extensive experiments across diverse long-horizon benchmarks and multiple frozen backbones demonstrate that TS-Memory consistently improves both point and probabilistic forecasting accuracy with constant-time inference overhead, enabling strictly retrieval-free deployment while preserving the robustness of retrieval-based adaptation.

% % Promising directions include continual memory updates under non-stationary drift while preserving retrieval-free inference, richer offline teachers via hybrid similarity and ensembles with utility-calibrated reliability learning, and leakage-safe alignment families that handle diverse non-stationarities such as scale, trend, and temporal warping. Further opportunities lie in input-adaptive fusion that adjusts memory reliance based on uncertainty, as well as scalable teacher compilation and extensions to broader tasks including imputation, anomaly detection, and decision-focused forecasting.
% Promising directions include continual memory updates under distribution drift while preserving retrieval-free inference, richer offline teachers with utility-calibrated reliability learning, and leakage-safe alignment for diverse non-stationarities. Further opportunities lie in input-adaptive fusion that adjusts memory reliance based on uncertainty, and extensions to broader tasks such as imputation, anomaly detection, and decision-focused forecasting(Appendix~\ref{sec:future_work}).

% We introduce \textbf{Parametric Memory Distillation}, a paradigm that compiles retrieval-conditioned knowledge offline into compact parametric memory, resolving the efficiency-adaptability dilemma in TSFM adaptation. Instantiated as \textbf{TS-Memory}, our approach distills quantile corrections from a privileged kNN teacher, equipping frozen backbones with domain-specific memory without runtime database dependencies. This achieves retrieval-level adaptability with constant-time inference and plug-and-play deployment. Experiments demonstrate that internalizing retrieval into neural memory eliminates latency while improving accuracy. Future work will explore continual memory updates and enhanced teacher reliability.
% We introduce \textbf{Parametric Memory Distillation}, a paradigm that compiles retrieval-conditioned knowledge offline into compact parametric memory, addressing the efficiency-adaptability dilemma in TSFM adaptation. Realized as \textbf{TS-Memory}, our approach distills quantile corrections from a privileged kNN teacher, equipping frozen backbones with domain-specific memory at constant-time inference and without runtime database dependencies. Experiments confirm that internalizing retrieval into neural memory removes search overhead while improving accuracy. Future directions include continual memory updates under evolving domains and richer teacher constructions with improved reliability estimation.

% We introduce Parametric Memory Distillation, a retrieval-to-memory paradigm that distills retrieval-conditioned predictive distributions into a lightweight, plug-and-play module for adapting frozen Time Series Foundation Models. Instantiated as TS-Memory, our approach learns confidence-aware quantile corrections from an offline, leakage-safe $k$NN teacher and enables retrieval-free inference through simple fusion with the frozen backbone. Experiments across multiple TSFM backbones and long-horizon benchmarks demonstrate consistent improvements in both point and probabilistic forecasting with negligible inference overhead. Future directions include continual memory updates under evolving domains and richer teacher constructions with improved reliability estimation.

We introduce \textbf{Parametric Memory Distillation}, a retrieval-to-memory paradigm that compiles retrieval-conditioned predictive distributions into a lightweight, plug-and-play module for adapting frozen Time Series Foundation Models. Realized as \textbf{TS-Memory}, our approach distills confidence-aware quantile corrections from an offline, leakage-safe $k$NN teacher and enables retrieval-free inference through simple fusion with the frozen backbone. Experiments across multiple TSFM backbones and long-horizon benchmarks demonstrate consistent improvements in both point and probabilistic forecasting with negligible inference overhead. Future directions include continual memory updates under evolving domains and richer teacher constructions with improved reliability estimation.

\section*{Limitations}

TS-Memory focuses on long-horizon probabilistic forecasting; extending it to imputation, anomaly detection, and decision-focused forecasting remains future work. 
Its shift alignment mainly handles additive channel-wise offsets, so temporal warping, frequency shifts, or abrupt regime changes may reduce reliability. 
When the retrieval corpus poorly matches the deployment domain, teacher targets can be noisy; the advantage gate, anchoring loss, and backbone fusion mitigate but cannot eliminate this risk. 
We use a validation-tuned fusion weight for each dataset--backbone pair at test time, leaving input-adaptive fusion for future study.

\section*{Acknowledgement}

We thank the reviewers for their valuable comments and efforts to improve this manuscript. This work is mainly supported by the Guangdong Basic and Applied Basic Research Foundation (No. 2025A1515011994). This work is also supported by the National Natural Science Foundation of China (No. 62402414), Guangdong Provincial Project 2025D03J0014, Guangzhou Municipal Science and Technology Project (No. 2023A03J0011), the Guangzhou Industrial Information and Intelligent Key Laboratory Project (No. 2024A03J0628), and Guangdong Provincial Key Lab of Integrated Communication, Sensing and Computation for Ubiquitous Internet of Things (No. 2023B1212010007). Raymond Chi-Wing Wong was supported by the fund PRP/004/25FX.